\chapter{Conclusion}\label{chap:conclusion}

This chapter summarizes the research findings, answers the research questions posed in Chapter~\ref{chap:intro}, discusses the contributions and limitations of this work, and suggests directions for future research.

\section{Summary of Research}

This thesis explored the applicability of Kubernetes as a platform for implementing business logic using Custom Resource Definitions (CRDs) and the Operator pattern. Two functionally equivalent billing systems were designed, implemented, and evaluated:

\begin{itemize}
    \item \textbf{backend-billing}: A traditional Go backend with REST API and PostgreSQL database
    \item \textbf{kube-billing}: A Kubernetes Operator with CRDs and reconcile-based control loop
\end{itemize}

Both systems implement identical business logic (billing plan management, subscription lifecycle, recurring billing, expiration handling) but differ fundamentally in architecture, state management, and execution model.

\section{Answers to Research Questions}

\subsection{RQ1: How can business logic be effectively modeled and implemented using Kubernetes CRDs and the Operator pattern?}

\textbf{Answer}: Business logic can be modeled using Kubernetes CRDs by:

\begin{enumerate}
    \item \textbf{Defining domain entities as CRDs}: BillingPlan and Subscription resources extend the Kubernetes API with domain-specific semantics
    \item \textbf{Using spec/status separation}: Desired state (spec) is declared by users; actual state (status) is managed by the controller
    \item \textbf{Implementing reconcile loop}: The Operator continuously compares spec vs status and applies changes to converge toward desired state
    \item \textbf{Leveraging Kubernetes primitives}: Finalizers for cleanup, Conditions for state reporting, Events for audit trail
    \item \textbf{Using admission control}: CRD validation markers enforce business rules at API boundary
\end{enumerate}

The key insight is that business logic becomes a \textit{control problem} rather than a \textit{computation problem}: the Operator acts as an automated control system that maintains system state according to business rules.

\subsection{RQ2: What are the performance characteristics of CRD/Operator-based implementation compared to traditional REST API architectures?}

\textbf{Answer}: The experimental evaluation (Chapter~\ref{chap:eval}) demonstrates significant performance differences:

\begin{table}[h]
\centering
\caption{Performance Comparison Summary}
\small\begin{tabularx}{\textwidth}{@{}lXX@{}}
\toprule
\textbf{Metric} & \textbf{backend-billing} & \textbf{kube-billing} \\ \midrule
Latency (p50, create subscription) & 10 ms & 100 ms \\
Throughput (requests/second) & $\sim$1,000 & $\sim$200 \\
CPU usage (idle) & 10m & 50m \\
Memory usage (idle) & 50 MB & 150 MB \\
Startup time & 2s & 10s \\ \bottomrule
\end{tabularx}
\end{table}

The traditional backend is \textbf{10--20$\times$ faster} in latency and \textbf{5$\times$ higher} in throughput. The Kubernetes API server and etcd introduce overhead from admission control, distributed consensus, and watch event propagation.

\textbf{Conclusion}: For performance-critical applications, the traditional approach is superior. The Operator pattern trades raw performance for operational benefits.

\subsection{RQ3: What are the operational implications and observability capabilities of using Kubernetes as a business logic platform?}

\textbf{Answer}: The Kubernetes approach provides significant operational advantages:

\begin{itemize}
    \item \textbf{Automatic scaling}: HorizontalPodAutoscaler adjusts replicas based on CPU/memory utilization
    \item \textbf{Self-healing}: Failed pods are automatically restarted by Kubernetes
    \item \textbf{Declarative deployment}: Kustomize/Helm enable version-controlled, reproducible deployments
    \item \textbf{Built-in observability}:
    \begin{itemize}
        \item Kubernetes Events provide audit trail without custom logging
        \item Conditions (metav1.Condition) standardize state reporting
        \item Reconciliation metrics expose control loop behavior
    \end{itemize}
    \item \textbf{GitOps compatibility}: CRD manifests integrate with ArgoCD/Flux workflows
\end{itemize}

However, these benefits come with operational complexity:
\begin{itemize}
    \item Kubernetes cluster management required
    \item etcd backup and disaster recovery planning
    \item RBAC configuration for least-privilege access
    \item Monitoring of control plane health
\end{itemize}

\subsection{RQ4: Under what circumstances is the CRD/Operator approach advantageous, and what are its limitations?}

\textbf{Answer}: The CRD/Operator approach is advantageous when:

\begin{itemize}
    \item \textbf{Application is already on Kubernetes} --- no additional infrastructure investment
    \item \textbf{High availability is critical} --- self-healing and HPA provide resilience
    \item \textbf{Business logic is state machine-like} --- reconcile pattern fits naturally
    \item \textbf{Audit trail is required} --- Kubernetes Events provide compliance logging
    \item \textbf{GitOps workflow is desired} --- declarative resources enable PR-based changes
\end{itemize}

The approach is \textbf{not recommended} when:

\begin{itemize}
    \item \textbf{Performance is the top priority} --- 10--20$\times$ latency overhead
    \item \textbf{Team lacks Kubernetes expertise} --- steep learning curve
    \item \textbf{Application must run outside Kubernetes} --- portability limitation
    \item \textbf{Complex queries are needed} --- Kubernetes API is not designed for ad-hoc queries
    \item \textbf{Time-to-market is critical} --- slower development cycle
\end{itemize}

\section{Contributions}

This thesis makes the following contributions:

\begin{enumerate}
    \item \textbf{Two production-ready implementations}:
    \begin{itemize}
        \item backend-billing: Go REST API with PostgreSQL, 86.2\% test coverage, CI/CD pipeline
        \item kube-billing: Kubernetes Operator with CRDs, HPA, PDB, production manifests
    \end{itemize}
    
    \item \textbf{Systematic architectural comparison}:
    \begin{itemize}
        \item Code metrics (LOC, complexity, dependencies, test coverage)
        \item Performance benchmarks (latency, throughput, resource usage)
        \item Operational characteristics (scaling, fault tolerance, observability)
    \end{itemize}
    
    \item \textbf{Decision framework}: Criteria for choosing between traditional and Operator approaches based on organizational priorities
    
    \item \textbf{Reference implementation}: Open-source codebases serving as templates for similar projects
\end{enumerate}

\section{Limitations}

Several limitations affect this research:

\begin{itemize}
    \item \textbf{Local environment}: Experiments conducted on Minikube/Docker Desktop; production clusters may show different characteristics
    
    \item \textbf{Simulated payments}: No integration with external payment gateway (Stripe, PayPal); real-world latency and failure modes not captured
    
    \item \textbf{Single workload type}: Billing domain may not represent all business logic patterns (e.g., real-time processing, batch jobs)
    
    \item \textbf{Short test duration}: No long-term effects (etcd compaction, database bloat, metric cardinality explosion)
    
    \item \textbf{Incomplete finalizer logic}: Pro-rated refund calculation not implemented in kube-billing
    
    \item \textbf{No multi-tenant testing}: Evaluation focused on single-tenant scenarios; multi-tenant isolation not assessed
\end{itemize}

\section{Future Work}

Several directions for future research emerge from this work:

\subsection{Technical Improvements}

\begin{itemize}
    \item \textbf{Payment gateway integration}: Implement Stripe/PayPal integration for realistic payment latency and failure handling
    
    \item \textbf{Refund logic}: Complete finalizer implementation with pro-rated refund calculation
    
    \item \textbf{Multi-tenant support}: Evaluate namespace-based isolation and resource quotas
    
    \item \textbf{Webhook integration}: Add ValidationWebhook and MutationWebhook for advanced admission control
    
    \item \textbf{Backup/restore}: Implement etcd backup strategies and disaster recovery procedures
\end{itemize}

\subsection{Research Extensions}

\begin{itemize}
    \item \textbf{Alternative domains}: Apply Operator pattern to other business domains (inventory management, order processing, user provisioning)
    
    \item \textbf{Hybrid architectures}: Explore combinations (e.g., REST API for user-facing operations, Operator for background processing)
    
    \item \textbf{Cost analysis}: Compare cloud infrastructure costs (managed Kubernetes vs managed database + compute)
    
    \item \textbf{Developer productivity}: Measure development velocity, onboarding time, and maintenance effort
    
    \item \textbf{Security analysis}: Evaluate RBAC configurations, secret management, and network policy effectiveness
\end{itemize}

\subsection{Production Deployment}

\begin{itemize}
    \item \textbf{Helm chart}: Package kube-billing as a Helm chart for easier deployment
    
    \item \textbf{Operator Lifecycle Manager}: Publish to OperatorHub for community distribution
    
    \item \textbf{Performance tuning}: Optimize etcd configuration, reconcile concurrency, and cache settings
    
    \item \textbf{Monitoring dashboards}: Create Grafana dashboards for business metrics visualization
\end{itemize}

\section{Final Remarks}

This thesis has demonstrated that Kubernetes can serve as a platform for business logic implementation through CRDs and the Operator pattern. The approach is technically feasible and provides operational benefits (automatic scaling, self-healing, built-in observability) but comes with significant performance overhead (10--20$\times$ latency, 5$\times$ lower throughput) and development complexity.

The choice between traditional and Operator-based architectures is not binary but contextual. Organizations should evaluate their specific requirements:

\begin{itemize}
    \item If \textbf{performance, simplicity, and portability} are priorities $\rightarrow$ traditional backend
    \item If \textbf{operational automation, GitOps, and cloud-native integration} are priorities $\rightarrow$ Kubernetes Operator
\end{itemize}

The research contributes empirical evidence to the cloud-native architecture discourse, moving beyond anecdotal claims to data-driven decision making. Both approaches have merit; the optimal choice depends on organizational context, team expertise, and business requirements.

\section{Publications and Artifacts}

The following artifacts are available as open-source:

\begin{itemize}
    \item \textbf{backend-billing}: \url{https://github.com/cuprum-acid/backend-billing}
    \item \textbf{kube-billing}: \url{https://github.com/cuprum-acid/kube-billing}
    \item \textbf{Thesis source}: \url{https://github.com/cuprum-acid/thesis}
\end{itemize}
